% GENERATED FILE. Edit canonical lessons, then run npm run generate:books.
% canonical-source-hash: fnv1a64:30f2157cd2123f47
% canonical-lessons: ZH-C325-jiangshu, ZH-C325-qushi, ZH-W325-shi-world, ZH-C325-houdai, ZH-C325-shidai, ZH-C325-xiangnian

\chapter{Gone but Remembered}
\label{ch:zh-gone-but-remembered}
\hlchaptermodality{325}

\begin{chapteropening}
\textbf{By the end of this chapter:} \emph{I can give an account of the past, say someone has died, and talk about descendants, generations and missing someone.}

Complete the last lesson of chapter 325: I can give an account of the past, say someone has died, and talk about descendants, generations and missing someone.
\end{chapteropening}

\section[jiǎngshù]{jiǎngshù — to tell about; to give an account of}

\label{lesson:ZH-C325-jiangshu}



\begin{quote}
Before the new one: say the Mandarin for an age in history, then the Mandarin for an impression left on you.
\end{quote}

\subsection*{You'll want to know: jiǎngshù}
\emph{jiǎngshù} — \textquotedblleft{}to tell about; to give an account of\textquotedblright{}.

\begin{quote}
\emph{tā gěi wǒmen jiǎngshù le nàge shíqī de shēnghuó} — He told us what life was like in that period.
\end{quote}

Its characters come later; for now it is a word for the ear.

\begin{cousinweb}
\emph{jiǎngshù} is built from \emph{jiǎng} and \emph{shù}, both inherited from Middle Chinese, \emph{kaewngX} and \emph{zywit} (Guangyun, in Baxter's transcription), which glosses the first 'to tell' and 'to discuss' and the second 'to set forth in writing'.
\end{cousinweb}

\begin{grammarlens}[title={What you've built}]
One more word for telling what happened.
\end{grammarlens}

\subsection*{Your turn}
\begin{itemize}
\raggedright
  \item \emph{Say it:} \emph{jiǎngshù}
  \item \emph{Say it:} \emph{jiǎngshù}, once more
  \item \emph{Say it:} \emph{tā gěi wǒmen jiǎngshù le nàge shíqī de shēnghuó} — He told us what life was like in that period.
  \item \emph{Recall:} say the Mandarin for an age in history, then the Mandarin for an impression left on you, then say \emph{jiǎngshù} again
\end{itemize}

\begin{tcolorbox}[breakable,colback=teal!4,colframe=teal!35!black,title={Before you move on}]
What does jiǎngshù mean? (\textquotedblleft{}To tell about; to give an account of\textquotedblright{}.) Say it once more.
\end{tcolorbox}
\section[\texorpdfstring{\zh{去世} (qùshì)}{去世 (qùshì)}]{\zh{去世} (qùshì) — to pass away}

\label{lesson:ZH-C325-qushi}



\begin{quote}
Before the new one: say the Mandarin for an impression left on you, then the Mandarin for to give an account of.
\end{quote}

\subsection*{You'll want to know: \zh{去世}}
\textbf{\zh{去世}} — \emph{qùshì} — \textquotedblleft{}to pass away\textquotedblright{}.

\begin{quote}
\textbf{\zh{他的外公去年去世了}} \emph{tā de wàigōng qùnián qùshì le} — His grandfather died last year.
\end{quote}

The new character, in pieces: \zh{世} — five strokes: a long bar; two verticals through it; a short bar joining their feet; then a stroke down the left that turns right along the bottom. The next lesson puts it in your hand.

\begin{cousinweb}
\textbf{\zh{去世}} (\emph{qùshì}) is built from \zh{去} 'to leave' and \zh{世} 'the world', both inherited from Middle Chinese, \emph{khjoH} and \emph{syejH} (Guangyun, in Baxter's transcription), which glosses the first 'to leave' and the second 'a generation'. The sense 'the world' is not given there.
\end{cousinweb}

\begin{grammarlens}[title={What you've built}]
One more word for telling what happened.
\end{grammarlens}

\subsection*{Your turn}
\begin{itemize}
\raggedright
  \item \emph{Say it:} \emph{qùshì}
  \item \emph{Say it:} \emph{qùshì}, once more
  \item \emph{Say it:} \emph{tā de wàigōng qùnián qùshì le} — His grandfather died last year.
  \item \emph{Recall:} say the Mandarin for an impression left on you, then the Mandarin for to give an account of, then say \emph{qùshì} again
\end{itemize}

\begin{tcolorbox}[breakable,colback=teal!4,colframe=teal!35!black,title={Before you move on}]
What does \zh{去世} mean? (\textquotedblleft{}To pass away\textquotedblright{}.) Say it once more.
\end{tcolorbox}
\section[\texorpdfstring{\zh{世} (shì)}{世 (shì)}]{\zh{世} — a long bar and two verticals}

\label{lesson:ZH-W325-shi-world}



\begin{quote}
The word first, then the character inside it.

\begin{itemize}
\raggedright
  \item \emph{Recall:} say \emph{qùshì} — \textbf{R1}, one lesson back, before it decays
\end{itemize}

\begin{quote}
\zh{十} \zh{去}
\end{quote}

\emph{Read it:} each one aloud, with nothing said first — \emph{shí}, \emph{qù}
\end{quote}

\subsection*{Script — a long bar and two verticals}
\hlblockfigure{\detokenize{figures/ZH-W325-shi-world-filmstrip.pdf}}{How \zh{世} is written, stroke by stroke}

\begin{quote}
\zh{世}
\end{quote}

\textbf{One piece, five strokes.}

\textbf{First the long bar,} from left to right.

\textbf{Then two verticals through it,} the left one first, each longer above the bar than below it.

\textbf{Then a short bar} joining the feet of the two verticals.

\textbf{Last, the turning stroke:} down the left side, then right along the bottom, under everything.

1 + 2 + 1 + 1 = five.

\subsection*{Your turn}
\begin{itemize}
\raggedright
  \item \emph{Write it:} \zh{世}, all five, with the model in front of you
  \item \emph{Cover:} write it again from memory, and count the pieces, not the strokes
\end{itemize}

\begin{tcolorbox}[breakable,colback=teal!4,colframe=teal!35!black,title={Before you move on}]
How many strokes in \zh{世}? (\textbf{Five}.) What comes after the long bar? (\textbf{Two verticals through it}.) What is the last stroke? (\textbf{One that runs down the left and turns right along the bottom}.)

Source: \href{https://raw.githubusercontent.com/chanind/hanzi-writer-data/68d10a4b21150cae5e1ebbd223eed289cf32d90c/data/\%E4\%B8\%96.json}{Hanzi Writer Data \zh{世}, pinned stroke paths and medians}.
\end{tcolorbox}
\section[\texorpdfstring{\zh{后代} (hòudài)}{后代 (hòudài)}]{\zh{后代} (hòudài) — descendants; later generations}

\label{lesson:ZH-C325-houdai}



\begin{quote}
Before the new one: say the Mandarin for to give an account of, then the Mandarin for to pass away.
\end{quote}

\subsection*{You'll want to know: \zh{后代}}
\textbf{\zh{后代}} — \emph{hòudài} — \textquotedblleft{}descendants; later generations\textquotedblright{}.

\begin{quote}
\textbf{\zh{他的后代现在都在国外}} \emph{tā de hòudài xiànzài dōu zài guówài} — His descendants are all abroad now.
\end{quote}

\begin{cousinweb}
\textbf{\zh{后代}} (\emph{hòudài}) is built from \zh{后} 'after' and \zh{代} 'a generation', both inherited from Middle Chinese, \emph{huwX} and \emph{dojH} (Guangyun, in Baxter's transcription), which glosses the first 'before and after' and the second 'to take turns'.
\end{cousinweb}

\begin{grammarlens}[title={What you've built}]
One more word for telling what happened.
\end{grammarlens}

\subsection*{Your turn}
\begin{itemize}
\raggedright
  \item \emph{Say it:} \emph{hòudài}
  \item \emph{Say it:} \emph{hòudài}, once more
  \item \emph{Say it:} \emph{tā de hòudài xiànzài dōu zài guówài} — His descendants are all abroad now.
  \item \emph{Recall:} say the Mandarin for to give an account of, then the Mandarin for to pass away, then say \emph{hòudài} again
\end{itemize}

\begin{tcolorbox}[breakable,colback=teal!4,colframe=teal!35!black,title={Before you move on}]
What does \zh{后代} mean? (\textquotedblleft{}Descendants; later generations\textquotedblright{}.) Say it once more.
\end{tcolorbox}
\section[\texorpdfstring{\zh{世代} (shìdài)}{世代 (shìdài)}]{\zh{世代} (shìdài) — for generations}

\label{lesson:ZH-C325-shidai}



\begin{quote}
Before the new one: say the Mandarin for to pass away, then the Mandarin for one's descendants.
\end{quote}

\subsection*{You'll want to know: \zh{世代}}
\textbf{\zh{世代}} — \emph{shìdài} — \textquotedblleft{}for generations\textquotedblright{}.

\begin{quote}
\textbf{\zh{他们家世代都是医生}} \emph{tāmen jiā shìdài dōu shì yīshēng} — Their family have been doctors for generations.
\end{quote}

\begin{cousinweb}
\textbf{\zh{世代}} (\emph{shìdài}) is built from \zh{世} 'a generation' and \zh{代} 'an age', both inherited from Middle Chinese, \emph{syejH} and \emph{dojH} (Guangyun, in Baxter's transcription), which glosses the first with the second: a generation.
\end{cousinweb}

\begin{grammarlens}[title={What you've built}]
One more word for telling what happened.
\end{grammarlens}

\subsection*{Your turn}
\begin{itemize}
\raggedright
  \item \emph{Say it:} \emph{shìdài}
  \item \emph{Say it:} \emph{shìdài}, once more
  \item \emph{Say it:} \emph{tāmen jiā shìdài dōu shì yīshēng} — Their family have been doctors for generations.
  \item \emph{Recall:} say the Mandarin for to pass away, then the Mandarin for one's descendants, then say \emph{shìdài} again
\end{itemize}

\begin{tcolorbox}[breakable,colback=teal!4,colframe=teal!35!black,title={Before you move on}]
What does \zh{世代} mean? (\textquotedblleft{}For generations\textquotedblright{}.) Say it once more.
\end{tcolorbox}
\section[xiǎngniàn]{xiǎngniàn — to miss someone}

\label{lesson:ZH-C325-xiangnian}



\begin{quote}
Before the new one: say the Mandarin for one's descendants, then the Mandarin for for generations on end.
\end{quote}

\subsection*{You'll want to know: xiǎngniàn}
\emph{xiǎngniàn} — \textquotedblleft{}to miss someone\textquotedblright{}.

\begin{quote}
\emph{wàigōng qùshì yǐhòu wǒ hěn xiǎngniàn tā} — After my grandpa died I missed him very much.
\end{quote}

Its characters come later; for now it is a word for the ear.

\begin{cousinweb}
\emph{xiǎngniàn} is built from \zh{想} 'to think' and \emph{niàn} 'to think of', both inherited from Middle Chinese, \emph{sjangX} and \emph{nemH} (Guangyun, in Baxter's transcription), which glosses both 'to think'.
\end{cousinweb}

\begin{grammarlens}[title={What you've built}]
One more word for telling what happened.
\end{grammarlens}

\subsection*{Your turn}
\begin{itemize}
\raggedright
  \item \emph{Say it:} \emph{xiǎngniàn}
  \item \emph{Say it:} \emph{xiǎngniàn}, once more
  \item \emph{Say it:} \emph{wàigōng qùshì yǐhòu wǒ hěn xiǎngniàn tā} — After my grandpa died I missed him very much.
  \item \emph{Recall:} say the Mandarin for one's descendants, then the Mandarin for for generations on end, then say \emph{xiǎngniàn} again
\end{itemize}

\begin{tcolorbox}[breakable,colback=teal!4,colframe=teal!35!black,title={Before you move on}]
What does xiǎngniàn mean? (\textquotedblleft{}To miss someone\textquotedblright{}.) Say it once more.
\end{tcolorbox}
